\section{第三部分：为 Lean 构建 Hammer——LeanHammer}

\subsection{Hammer 管道架构}

\begin{importantbox}{LeanHammer 组件}
\begin{enumerate}
    \item \textbf{前提选择}（Premise Selection）：神经检索模型嵌入目标和候选前提，用余弦相似度选择最相关的定理/定义/引理
    \item \textbf{翻译层}（lean-auto）：将 Lean 的依赖类型论翻译为自动证明器的逻辑
    \item \textbf{自动定理证明器}（Zipperposition 等）
    \item \textbf{证明重构}（Duper）：将自动证明器的输出翻译回 Lean 证明
    \item \textbf{树搜索}（AESOP）：组织整个搜索过程
\end{enumerate}
\end{importantbox}

用户只需在证明中调用一个 \texttt{hammer} 命令，系统自动完成前提选择、翻译、自动证明和重构。

\subsection{实验结果}

\begin{itemize}
    \item LeanHammer 的神经前提选择器（仅 1 亿参数）显著优于先前工作
    \item 与符号前提选择器（Moogle）互补，取并集后性能进一步提升
    \item 性能正在接近使用人类证明中的真实前提时的理论上限
\end{itemize}

\subsection{本章小结}

DSP 和 LeanHammer 代表了两种互补的方法：前者通过高级草图引导证明结构，后者提供强大的低级自动化填补工具。

